\chapter{Avaliação e Discussão}
\label{cap:avaliacao}

Este capítulo confronta o que foi construído com o problema de pesquisa e com os objetivos declarados na introdução. A avaliação combina uma comparação estruturada entre o portal e o \textit{dashboard} com a discussão das limitações remanescentes e das implicações para a transparência e o controle social. Não houve teste de usabilidade com usuários finais, a avaliação aqui reportada é autoral, apoiada na reprodução repetida das mesmas perguntas nos dois ambientes.

\section{Avaliação da Usabilidade}

O critério de comparação adotado é funcional: "dada uma pergunta de fiscalização, o sistema permite respondê-la com quantos passos, em quanto tempo e com que grau de interpretação técnica intermediária?" Não se mede a estética da interface nem a satisfação subjetiva, o que se mede é a capacidade analítica relativa. A Tabela~\ref{tab:comparacao-portal-dashboard} resume seis tarefas representativas.

\begin{table}[htb]
\caption{Comparação portal PJF \textit{versus} \textit{dashboard} em tarefas de fiscalização}
\label{tab:comparacao-portal-dashboard}
\centering
\small
\renewcommand{\arraystretch}{1.5}
\begin{tabular*}{\textwidth}{@{\extracolsep{\fill}}p{4.2cm}p{4.8cm}p{4.8cm}@{}}
\hline
\textbf{Pergunta} & \textbf{Portal da Transparência} & \textbf{\textit{Dashboard}} \\
\hline
Saldo entre arrecadação e pagamentos no ano & Baixar planilhas de receita e despesa; alinhar competência; somar colunas manualmente & Página Visão Geral Fiscal; filtrar ano \\

Evolução mensal empenhado/liquidado/pago & Vários arquivos mensais; colar e consolidar & Página Despesas; gráfico único \\

Pagamentos por função & Filtrar planilha por código funcional; somar; decodificar código & Página Execução Orçamentária; filtrar anos; ler tabela com rótulos LOA \\

Participação dos maiores fornecedores & Exportar; ordenar; calcular \% no Excel & Página Fornecedores; indicador "\% nos 5 maiores" \\

Previsão \textit{vs.} arrecadação por natureza & Reconciliar comparativo e previsão; interpretar códigos & Página Receitas; tabela com \% de realização \\

Significado de "liquidação" & Buscar glossário externo ou norma & Glossário integrado + \textit{tooltips} \\
\hline
\end{tabular*}
\end{table}

Em todas as tarefas listadas, o portal contém a informação bruta, mas externaliza ao usuário o trabalho de integração. O \textit{dashboard} não cria novos dados: reorganiza o que já foi publicado e define convenções que, uma vez implementadas no \textit{warehouse}, não precisam ser refeitas a cada consulta. O ganho de usabilidade é, portanto, herdado da camada analítica descrita no Capítulo~\ref{cap:dw}, não de um truque de visualização isolado.

\section{Limitações do Trabalho}

As limitações do \textit{warehouse} (Capítulo~\ref{cap:dw}) propagam-se ao \textit{dashboard}. Ausência de dotação autorizada impede indicadores de execução sobre o previsto na LOA, métrica padrão em relatórios de controle interno. O grão por empenho existe no \textit{warehouse}, e \texttt{dim\_empenho} guarda a descrição e os metadados de licitação. O que falta é o detalhe abaixo do empenho. O pagamento aparece apenas como total mensal, sem data própria, o que dificulta a medição de prazos até a quitação. Itens, quantidades e preços de cada empenho não são publicados, e a interface ainda não usa esses metadados. Rótulos LOA com \textit{fallback} em cerca de 14\,\% dos pares função/subfunção presentes na execução comprometem a legibilidade pontual.

Limitações próprias da implementação do painel acrescentam-se às da fonte. O painel está publicado no \textit{Streamlit Community Cloud}\footnote{\url{https://juiz-de-fora-transparencia.streamlit.app/}}, lendo um DuckDB em arquivo empacotado na aplicação. O orquestrador Dagster, porém, não roda de forma agendada e, por isso, a atualização exige reexecutar manualmente o \textit{pipeline} e republicar o banco de dados. Os dados exibidos refletem, portanto, a última materialização, não necessariamente o portal em tempo real. Por fim, o cruzamento com licitações e contratos, citado na introdução como uma lacuna dos portais, permanece fora do escopo, embora conste como trabalho futuro.

\section{Implicações para Transparência e Controle Social}

Retomando as três lacunas da introdução, o sistema construído endereça cada uma delas em graus distintos.

A lacuna técnica é a mais diretamente atacada. Integrar receitas e despesas, conciliar planilhas heterogêneas, normalizar naturezas com o ementário STN e expor séries temporais filtráveis eliminam o gargalo de consolidação manual que a literatura associa aos portais municipais \cite{marco2022transparencia,melo2019proposta}. O cidadão deixa de depender de uma planilha intermediária montada por um especialista ao acessar o \textit{dashboard} em uma URL pública, sem instalar o ambiente localmente.

A lacuna social é parcialmente mitigada. Glossário, linguagem de "visão cidadã" e agrupamento das naturezas de receita por origem aproximam os termos do MCASP do vocabulário público. Permanecem os códigos nos filtros e a distância entre a classificação contábil e a experiência vivida: o painel mostra quanto foi pago em "saúde" na classificação funcional, não quanto foi investido em um equipamento específico no bairro do usuário. Transparência agregada educa sobre prioridades macro; não substitui prestação de contas situada.

A lacuna pedagógica encontra uma resposta incompleta. Freire \cite{freire1967liberdade} exige que o sujeito formule perguntas e construa interpretações, não apenas decodifique tabelas. Os \textit{expanders} orientadores, a estrutura por perguntas e a seção analítica do Capítulo~\ref{cap:resultados} empurram nessa direção. Kimball \cite{kimball2013data} fornece a estrutura orientada a assunto que torna a pergunta operacionalizável em SQL e, em seguida, em um gráfico. O encontro das duas tradições se materializa quando o usuário filtra por saúde, compara anos e questiona a concentração de fornecedores: gesto de "ler o mundo" orçamentário com ferramenta, não só para consumir um relatório passivo.

\section{Replicabilidade e Escalabilidade}

O modelo é replicável para outros municípios com portais que publiquem planilhas de execução em formato semelhante, mas a replicação não é uma cópia literal. Cada cidade impõe \textit{parsers} próprios. A arquitetura em camadas, extração com Dagster, \textit{staging} com dbt, \textit{warehouse} dimensional, consumo com Streamlit, permanece estável, o que muda são \textit{assets} de extração e dimensões auxiliares. Adaptar para vinte municípios mineiros exigiria, na prática, vinte conjuntos de regras de \textit{parsing} documentadas, preferencialmente com testes automatizados e política de manutenção quando os portais mudarem.

Escalar também é decisão institucional. Manter o \textit{pipeline} atualizado e o painel no ar exige uma rotina de \textit{backfill} e um responsável técnico, papel que, hoje, recai sobre o autor do TCC. A instância pública no \textit{Streamlit Community Cloud} reduz a barreira de acesso, mas não substitui a governança de dados nem o patrocínio institucional de longo prazo. Iniciativas \textit{civic tech} mostram impacto quando focadas \cite{marco2022transparencia}; sistemas institucionais duram mais, porém conservam recortes. Um meio-termo viável seria um observatório civil ou uma universidade local assumir a manutenção do \textit{pipeline} e da instância hospedada, com código aberto e documentação de reprodução já presentes no repositório do trabalho.

Para Juiz de Fora especificamente, o \textit{dashboard} já está disponível em uma URL pública. O próximo passo de maior retorno social seria automatizar a atualização do \textit{pipeline} e validar as tarefas da Tabela~\ref{tab:comparacao-portal-dashboard} com usuários reais. Sem essa cadência e sem evidência de uso pelo cidadão, a contribuição permanece uma ferramenta acessível, porém ainda pouco consolidada no debate democrático local.
